% =====================================================================
%  ch03a.tex —— 第 3 章 3.1–3.2（原书 p.544–549 前）
% =====================================================================
% [原文 p.544]
\chapter{导出范畴之间的全忠实函子}

\section{Postnikov 图及其卷积}

\takeaway{本节在三角范畴中引入 Postnikov 系统：把有界复形 $X^{\cdot}$ 的微分逐级换成映射锥，串成首尾相接的特异三角，末项 $Y^{0}$ 称为 $X^{\cdot}$ 的卷积。目标是给出卷积存在、且在同构意义下唯一的充分条件，即消没条件 (25)–(28)。}

% [Tutor Note: 原书节标题与正文接排（“3.1. Postnikov diagrams and their convolutions. In this section…”），本译本按全稿体例改为独立节标题，并加节首重点框。]

本节考虑三角范畴中的 Postnikov 图，并给出使一个 Postnikov 图容许卷积、且该卷积在同构意义下唯一确定的条件。

设 $c<0$，而 $X^{\cdot}=\{X^{c}\xrightarrow{d^{c}}X^{c+1}\xrightarrow{d^{c+1}}\cdots\to X^{0}\}$ 是三角范畴 $\cD$ 中对象的有界复形。这意思是说，所有复合 $d^{i+1}\circ d^{i}$ 都为零。

按定义，与 $X^{\cdot}$ 相伴的一个\emph{左 Postnikov 系统}是如下形式的图：
\begin{center}
\begin{tikzcd}[row sep=1.9em,column sep=1.45em,ampersand replacement=\&]
X^{c} \arrow[rr,"d^{c}"] \arrow[ddd,"i_{c}=\id"'] \& {} \&
X^{c+1} \arrow[rr,"d^{c+1}"] \arrow[ddd,"i_{c+1}"'] \& {} \&
X^{c+2} \arrow[ddd,"i_{c+2}"'] \& \cdots \& {} \& {} \&
X^{0} \arrow[ddd,"i_{0}"'] \\
{} \& \circlearrowleft \& {} \& \circlearrowleft \& {} \& {} \& {} \& {} \& {} \\
{} \& \star \& {} \& \star \& {} \& {} \& {} \& \star \& {} \\
Y^{c}=X^{c} \arrow[uuurr,"j_{c}"] \& {} \&
Y^{c+1} \arrow[ll,"{[1]}"'] \arrow[uuurr,"j_{c+1}"] \& {} \&
Y^{c+2} \arrow[ll,"{[1]}"'] \& \cdots \&
Y^{-1} \arrow[uuurr,"j_{-1}"] \& {} \&
Y^{0} \arrow[ll,"{[1]}"']
\end{tikzcd}
\end{center}
% [Tutor Note: 页图上斜箭头的落点在 200 dpi 下压着底行字缘，无法直读；此处按 $j_{k}\circ i_{k}=d^{k}$ 与 $Y^{c}=X^{c}$、$i_{c}=\id$ 的相容性判定为 $i_{k}:X^{k}\to Y^{k}$、$j_{k}:Y^{k}\to X^{k+1}$，与原文同段所述的交换三角一致。]
% [原文 p.545]
其中标以 $\star$ 的三角均为特异三角，标以 $\circlearrowleft$ 的均为交换的（也就是说，$j_{k}\circ i_{k}=d^{k}$）。若存在与 $X^{\cdot}$ 相伴的左 Postnikov 系统使 $E=Y^{0}$，则称对象 $E\in\Ob\cD$ 为复形 $X^{\cdot}$ 的一个\emph{左卷积}。用 $\Tot(X^{\cdot})$ 记复形 $X^{\cdot}$ 的所有卷积构成的类。显然，Postnikov 系统及其卷积在三角范畴之间的正合函子下保持。

\yizhu{Postnikov 系统是三角范畴中对“全复形”的替代：三角范畴没有核与余核，无法像在复形范畴里那样直接取 $\Tot$，这里的做法是逐级把相邻两项的微分换成映射锥，得到一串首尾相接的特异三角，末项 $Y^{0}$ 就是卷积。同一个复形可以有互不同构的卷积，$\Tot(X^{\cdot})$ 因此一般只是一个类，既可能含多个对象，也可能为空。〔译者按（上下文推断，无外部依据）：图中圆圈标出的是 $j_{k}\circ i_{k}=d^{k}$ 这一交换性，正是它保证逐级的锥与原微分相容。〕【来源：原书 §3.1，p.544–545】}

注意，类 $\Tot(X^{\cdot})$ 可能含有许多互不同构的对象，也可能是空的。下面我们将给出一个充分条件，使得类 $\Tot(X^{\cdot})$ 在同构意义下只含一个对象。下面的引理在 [3] 中证明。

\begin{lemma}{3.1.1}
设 $g$ 是对象 $Y$ 与 $Y'$ 之间的态射，而 $Y$ 与 $Y'$ 又分别落在特异三角中：
\diag{
X \arrow[r,"u"] \arrow[d,"f"'] \& Y \arrow[r,"v"] \arrow[d,"g"'] \& Z \arrow[r,"w"] \arrow[d,"h"'] \& X[1] \arrow[d,"{f[1]}"'] \\
X' \arrow[r,"u'"'] \& Y' \arrow[r,"v'"'] \& Z' \arrow[r,"w'"'] \& X'[1]
}
若 $v'gu=0$，则存在态射 $f:X\to X'$ 与 $h:Z\to Z'$，使得三元组 $(f,g,h)$ 是三角态射。

此外，若还假设 $\Hom(X[1],Z')=0$，则态射 $f$ 与 $h$（它们分别使图中第一个与第二个方块交换）由这些条件唯一确定。
\end{lemma}

下面我们证明两个引理，它们把上一个引理推广到 Postnikov 图。

\begin{lemma}{3.1.2}
设 $X^{\cdot}=\{X^{c}\xrightarrow{d^{c}}X^{c+1}\xrightarrow{d^{c+1}}\cdots\to X^{0}\}$ 是三角范畴 $\cD$ 中对象的有界复形。假设它对一切 $a<b$ 与 $i<0$ 满足
\begin{equation}
\Hom^{i}(X^{a},X^{b})=0.
\tag{25}
\end{equation}
则 $X^{\cdot}$ 的卷积存在，且所有卷积都（非典范地）同构。

此外，若还假设对某个卷积 $Y^{0}$（从而对任一卷积）、对一切 $a$ 与 $i<0$ 都有
\begin{equation}
\Hom^{i}(X^{a},Y^{0})=0,
\tag{26}
\end{equation}
则 $X^{\cdot}$ 的所有卷积都典范同构。
\end{lemma}

\begin{lemma}{3.1.3}
设 $X_{1}^{\cdot}$ 与 $X_{2}^{\cdot}$ 是满足条件 (25) 的有界复形，而 $(f_{c},\ldots,f_{0})$ 是这两个复形之间的态射：
\diag{
X_{1}^{c} \arrow[r,"d_{1}^{c}"] \arrow[d,"f_{c}"'] \& X_{1}^{c+1} \arrow[r] \arrow[d,"f_{c+1}"'] \& \cdots \arrow[r] \& X_{1}^{0} \arrow[d,"f_{0}"'] \\
X_{2}^{c} \arrow[r,"d_{2}^{c}"'] \& X_{2}^{c+1} \arrow[r] \& \cdots \arrow[r] \& X_{2}^{0}
}
% [原文 p.546]
假设对 $a<b$ 与 $i<0$ 有
\begin{equation}
\Hom^{i}(X_{1}^{a},X_{2}^{b})=0.
\tag{27}
\end{equation}
则对 $X_{1}^{\cdot}$ 的每个卷积 $Y_{1}^{0}$ 与 $X_{2}^{\cdot}$ 的每个卷积 $Y_{2}^{0}$，存在与态射 $f_{0}$ 交换的态射 $f:Y_{1}^{0}\to Y_{2}^{0}$。此外，若对一切 $a$ 与 $i<0$ 还有
\begin{equation}
\Hom^{i}(X_{1}^{a},Y_{2}^{0})=0,
\tag{28}
\end{equation}
则这个态射唯一确定。
\end{lemma}

\begin{proof}
我们对两个引理同时归纳证明，归纳的基础是引理 3.1.1。设 $Y^{c+1}$ 是态射 $d^{c}$ 的映射锥，
\begin{equation}
X^{c}\xrightarrow{d^{c}}X^{c+1}\xrightarrow{\alpha}Y^{c+1}\longrightarrow X^{c}[1].
\tag{29}
\end{equation}
按假设，$d^{c+1}\circ d^{c}=0$ 且 $\Hom(X^{c}[1],X^{c+2})=0$。因此存在唯一的态射 $\overline{d^{c+1}}:Y^{c+1}\to X^{c+2}$，使 $\overline{d^{c+1}}\circ\alpha=d^{c+1}$。考虑复合
\begin{equation*}
d^{c+2}\circ\overline{d^{c+1}}:Y^{c+1}\longrightarrow X^{c+3}.
\end{equation*}
已知 $d^{c+2}\circ\overline{d^{c+1}}\circ\alpha=d^{c+2}\circ d^{c+1}=0$；此外还有等式 $\Hom(X^{c}[1],X^{c+3})=0$。由此立即推出，复合 $d^{c+2}\circ\overline{d^{c+1}}$ 也为零。

% [Tutor Note: 文字层把 $\overline{d^{c+1}}$ 的上划线整个丢掉，抽成 $d^{c+1}$（例如“there is a unique morphism dc+1 : Y c+1 → Xc+2 such that dc+1 ◦ $\alpha$ = dc+1”），页图印作 $\overline{d^{c+1}}$；本译本按页图加上划线，以便与同段中的 $d^{c+1}$ 区分。]

考察特异三角 (29)，我们看出：对 $i<0$ 与 $b>c+1$，
\begin{equation*}
\Hom^{i}(Y^{c+1},X^{b})=0.
\end{equation*}
于是复形 $Y^{c+1}\to X^{c+2}\to\cdots\to X^{0}$ 也满足 (25)。由归纳假设，这个复形有卷积。因此 $X^{\cdot}$ 也有卷积，从而类 $\Tot(X^{\cdot})$ 非空。

下面我们证明：在条件 (27) 之下，每个复形态射都扩充为 Postnikov 系统的态射。考察态射 $d_{1}^{c}$ 与 $d_{2}^{c}$ 的映射锥 $Y_{1}^{c+1}$ 与 $Y_{2}^{c+1}$。存在态射 $g_{c+1}:Y_{1}^{c+1}\to Y_{2}^{c+1}$，把配对 $(f_{c},f_{c+1})$ 补成三角态射，
\diag{
X_{1}^{c} \arrow[r,"d_{1}^{c}"] \arrow[d,"f_{c}"'] \& X_{1}^{c+1} \arrow[r,"\alpha"] \arrow[d,"f_{c+1}"'] \& Y_{1}^{c+1} \arrow[r] \arrow[d,"g_{c+1}"'] \& X_{1}^{c}[1] \arrow[d,"{f_{c}[1]}"'] \\
X_{2}^{c} \arrow[r,"d_{2}^{c}"'] \& X_{2}^{c+1} \arrow[r,"\beta"'] \& Y_{2}^{c+1} \arrow[r] \& X_{2}^{c}[1]
}
如上文已经证明的，对 $i=1,2$ 存在唯一确定的态射 $\overline{d_{i}^{c+1}}:Y_{i}^{c+1}\to X_{i}^{c+2}$。考虑下图：
\diag{
Y_{1}^{c+1} \arrow[r,"\overline{d_{1}^{c+1}}"] \arrow[d,"g_{c+1}"'] \& X_{1}^{c+2} \arrow[d,"f_{c+2}"'] \\
Y_{2}^{c+1} \arrow[r,"\overline{d_{2}^{c+1}}"'] \& X_{2}^{c+2}
}
% [原文 p.547]
我们证明这个方块交换。事实上，记差为 $h=f_{c+2}\circ\overline{d_{1}^{c+1}}-\overline{d_{2}^{c+1}}\circ g_{c+1}$。我们有等式 $h\circ\alpha=f_{c+2}\circ d_{1}^{c+1}-\overline{d_{2}^{c+1}}\circ f_{c+1}=0$。而按引理的假设，$\Hom(X_{1}^{c}[1],X_{2}^{c+2})=0$。这立即推出 $h=0$。

于是我们得到复形态射
\diag{
Y_{1}^{c+1} \arrow[r,"\overline{d_{1}^{c+1}}"] \arrow[d,"g_{c+1}"'] \& X_{1}^{c+2} \arrow[r] \arrow[d,"f_{c+2}"'] \& \cdots \arrow[r] \& X_{1}^{0} \arrow[d,"f_{0}"'] \\
Y_{2}^{c+1} \arrow[r,"\overline{d_{2}^{c+1}}"'] \& X_{2}^{c+2} \arrow[r] \& \cdots \arrow[r] \& X_{2}^{0}
}
这些复形满足条件 (25) 与 (27)。由归纳假设，这两个复形之间的态射扩充为 Postnikov 系统之间的态射。这样我们就得到与 $X_{1}^{\cdot}$ 与 $X_{2}^{\cdot}$ 相伴的 Postnikov 系统之间的一个态射。

此外，容易看出：若所有态射 $f_{i}$ 都是同构，则 Postnikov 系统之间的这个态射也是同构。因此，若条件 (25) 成立，$\Tot(X^{\cdot})$ 中的所有对象都同构。

最后，考察在 Postnikov 图中出现的特异三角之间的态射，
\diag{
Y_{1}^{-1} \arrow[r,"j_{-1}"] \arrow[d,"g_{-1}"'] \& X_{1}^{0} \arrow[r,"i_{0}"] \arrow[d,"f_{0}"'] \& Y_{1}^{0} \arrow[r] \arrow[d,"g_{0}"'] \& Y_{1}^{-1}[1] \arrow[d,"{g_{-1}[1]}"'] \\
Y_{2}^{-1} \arrow[r,"j_{-1}"'] \& X_{2}^{0} \arrow[r,"i_{0}"'] \& Y_{2}^{0} \arrow[r] \& Y_{2}^{-1}[1]
}
若复形 $X_{i}^{\cdot}$ 满足条件 (28)（也就是说，对 $i<0$ 与一切 $a$ 有 $\Hom^{i}(X_{1}^{a},Y_{2}^{0})=0$），则 $\Hom(Y_{1}^{-1}[1],Y_{2}^{0})=0$。由引理 3.1.1，态射 $g_{0}$ 唯一确定。引理的证明至此完成。
\end{proof}

\section{凝聚层导出范畴之间的全忠实函子}

\takeaway{本节给出本章主定理：光滑射影簇之间凝聚层导出范畴的正合函子，只要全忠实且有伴随，就由积上的一个对象按式 (30) 表示，且该对象在同构意义下唯一；等价因必有伴随而自动可表示。另证有界性引理 3.2.3，为 §3.3 的构造铺路。}

设 $X$ 与 $M$ 是某个域 $k$ 上的两个光滑完备簇。如前面一样，用 $\Db(X)$ 与 $\Db(M)$ 分别记 $X$ 与 $M$ 上凝聚层的有界导出范畴。前面已经证明，这两个范畴具有三角范畴的结构。

考虑积 $M\times X$，并用 $p$ 与 $\pi$ 分别记 $M\times X$ 到 $M$ 与到 $X$ 的投影：
\begin{equation*}
M\xleftarrow{\ p\ }M\times X\xrightarrow{\ \pi\ }X.
\end{equation*}
对每个对象 $\cE\in\Db(M\times X)$，我们按 (7) 定义了从 $\Db(M)$ 到 $\Db(X)$ 的正合函子 $\Phi_{\cE}$：
\begin{equation}
\Phi_{\cE}(\,\cdot\,):=\Rf\pi_{*}\bigl(\cE\otL p^{*}(\,\cdot\,)\bigr).
\tag{30}
\end{equation}
函子 $\Phi_{\cE}$ 有左伴随函子与右伴随函子 $\Phi^{*}_{\cE}$ 与 $\Phi^{!}_{\cE}$，它们分别由公式 (8) 给出：
\begin{equation*}
\Phi^{*}_{\cE}(\,\cdot\,)
=\Rf p_{*}\bigl(\cE^{\vee}\otL\pi^{*}(\omega_{X}[\dim X]\otimes(\,\cdot\,))\bigr),
\end{equation*}
\begin{equation*}
\Phi^{!}_{\cE}(\,\cdot\,)
=\omega_{M}[\dim M]\otimes\Rf p_{*}\bigl(\cE^{\vee}\otL(\,\cdot\,)\bigr),
\end{equation*}
其中 $\omega_{X}$、$\omega_{M}$ 是 $X$ 与 $M$ 的典范层，而 $\cE^{\vee}:=\Rb\mathcal{H}om(\cE,\cO_{M\times X})$。

\yizhu{形如 (30) 的函子即所谓\emph{积分变换}：核对象 $\cE$ 取自积 $M\times X$，函子先沿 $p$ 拉回、再与 $\cE$ 作导出张量积、最后沿 $\pi$ 前推。上一段末尾的两个公式给出它的左、右伴随，其中 $\cE^{\vee}$ 是 $\cE$ 的对偶复形。定理 3.2.1 断言：只要正合函子全忠实且有伴随，它就一定是某个核对象的积分变换，且核在同构意义下唯一——这正是第~3 章要证的“等价总可由积上的对象表示”。【来源：原书引言式 (1) 与 §3.2 式 (30)、p.547–548】}

% [Tutor Note: 式 (30) 页图印作 $\mathbf{R}\pi_{*}$（$\mathbf{R}$ 上没有点），而引言式 (1) 同一类函子印作 $\mathbf{R}^{\bullet}q_{*}$；本译本照录式 (30) 的写法。]

% [原文 p.548]
为了研究两个簇何时具有等价的凝聚层导出范畴这一问题，并描述它们的自等价群，最好能对所有正合函子都有显式公式。有一个猜想：它们全都可由积上的对象表示，也就是说，都具有 (30) 的形式。然而目前还不知道这个断言是否成立。不过，这个猜想的一个特殊情形是成立的。也就是说：若一个函子全忠实并且有伴随函子，则它可由积上的一个对象表示。本章致力于证明这一事实。更确切地说，本章的主定理如下。

\begin{theorem}{3.2.1}
设 $F$ 是从范畴 $\Db(M)$ 到范畴 $\Db(X)$ 的正合函子，其中 $M$ 与 $X$ 是光滑射影簇。假设 $F$ 全忠实，并且有右（或左）伴随函子。则存在对象 $\cE\in\Db(M\times X)$，使得 $F$ 同构于由 (30) 定义的函子 $\Phi_{\cE}$，且对象 $\cE$ 在同构意义下唯一确定。
\end{theorem}

由此立即推出，每个等价都可由积上的一个对象表示，因为每个等价都有伴随，而伴随与某个拟逆函子重合。

\begin{theorem}{3.2.2}
设 $M$ 与 $X$ 是两个光滑射影簇。假设正合函子 $F:\Db(M)\iso\Db(X)$ 是三角范畴的等价。则存在对象 $\cE\in\Db(M\times X)$，它在同构意义下唯一，使得 $F$ 同构于函子 $\Phi_{\cE}$。
\end{theorem}

这些结果使我们能够描述凝聚层导出范畴之间的所有等价，并回答两个不同簇何时具有等价的凝聚层导出范畴这一问题。

在开始证明这些定理之前，我们作一点说明。设 $F$ 是从 $\Db(M)$ 到 $\Db(X)$ 的正合函子。若存在，则分别用 $F^{*}$ 与 $F^{!}$ 记 $F$ 的左伴随函子与右伴随函子。若左伴随 $F^{*}$ 存在，则右伴随 $F^{!}$ 也存在，并由公式
\begin{equation*}
F^{!}=S_{M}\circ F^{*}\circ S_{X}^{-1}
\end{equation*}
定义，其中 $S_{X}$ 与 $S_{M}$ 是范畴 $\Db(X)$ 与 $\Db(M)$ 的 Serre 函子。这两个函子存在，分别等于 $(\,\cdot\,)\otimes\omega_{X}[\dim X]$ 与 $(\,\cdot\,)\otimes\omega_{M}[\dim M]$（见 (6)）。

\yizhu{Serre 函子把 Serre 对偶范畴化：在光滑射影簇 $X$ 的凝聚层有界导出范畴上，它由 $S_{X}=(\,\cdot\,)\otimes\omega_{X}[\dim X]$ 给出（原书所引式 (6) 即此）。有了它，左伴随 $F^{*}$ 一旦存在，右伴随便由 $F^{!}=S_{M}\circ F^{*}\circ S_{X}^{-1}$ 给出，无须另行构造——这正是“有左伴随”与“有右伴随”可互相推出的依据。【来源：原书 p.548 及所引式 (6)】}

设 $F$ 是从导出范畴 $\Db(\cA)$ 到导出范畴 $\Db(\cB)$ 的正合函子。若存在 $z\in\Z$ 与 $n\in\N$，使得对任一对象 $A\in\cA$，上同调 $H^{i}(F(A))$ 当 $i\notin[z,z+n]$ 时都消没，则称 $F$ 是\emph{有界的}。

\yizhu{有界性说的是：$F$ 把 $\Db(M)$ 中对象的上同调整体压进一个长度固定为 $n$ 的区间 $[z,z+n]$，两个端点都不随对象改变。引理 3.2.3 表明，只要 $F$ 有左伴随，这一性质自动成立：证明取极丰富线丛 $\cL$，用 $\cO(k)$ 的右消解与一个谱序列控制住 $G(\cL^{k})$（$k<0$）的上同调，再由 $\cL$ 的丰富性回到 $F$ 本身。【来源：原书 §3.2，p.548–549】}

\begin{lemma}{3.2.3}
设 $M$ 与 $X$ 是射影簇，且 $M$ 是光滑簇。若正合函子 $F:\Db(M)\to\Db(X)$ 有左伴随，则 $F$ 有界。
\end{lemma}

\begin{proof}
我们用 $G:\Db(X)\to\Db(M)$ 记 $F$ 的左伴随，并在 $X$ 上取极丰富的线丛 $\cL$。它给出嵌入 $i:X\into\PP^{N}$。对任一 $k<0$，$\PP^{N}$ 上的层 $\cO(k)$ 有以层 $\cO(j)$（$j=0,1,\ldots,N$）表示的右消解，其形式为
\begin{equation*}
\cO(k)\xrightarrow{\ \sim\ }\{V_{0}\otimes\cO\to V_{1}\otimes\cO(1)\to\cdots\to V_{N}\otimes\cO(N)\to 0\},
\end{equation*}
% [原文 p.549]
其中所有 $V_{j}$ 都是向量空间 [2]。把这个消解限制到 $X$，就得到层 $\cL^{k}$ 以层 $\cL^{j}$（$j=0,1,\ldots,N$）表示的消解。由于对任一 $j=0,1,\ldots,N$，对象 $G(\cL^{j})$ 的非零上同调都属于某个区间，故可找到整数 $z'$ 与正整数 $n'$，使得上同调 $H^{l}(G(\cL^{k}))$ 对一切 $k\leqslant 0$ 与一切 $l\notin[z',z'+n']$ 都消没。这立即由谱序列
\begin{equation*}
E_{1}^{p,q}=V_{p}\otimes H^{q}(G(\cL^{p}))\Longrightarrow H^{p+q}(G(\cL^{i}))
\end{equation*}
的存在性推出。

% [Tutor Note: 谱序列右端页图与文字层都印作 $G(\cL^{i})$，而同一句的结论写作 $G(\cL^{k})$；据上下文此处应为 $k$，但照录原文，不改。]

设 $A\in\Db(M)$ 是任一对象。由于 $\cL$ 丰富，可知：若对某个取定的 $j$ 与一切 $i\ll 0$ 都有 $\Hom^{j}(\cL^{i},F(A))=0$，则上同调 $H^{j}(F(A))$ 消没。按假设，$G$ 是 $F$ 的左伴随。因此
\begin{equation*}
\Hom^{j}(\cL^{i},F(A))\cong\Hom^{j}(G(\cL^{i}),A).
\end{equation*}

现在考虑 $M$ 上的层 $\cF$。由于对一切 $i<0$，对象 $G(\cL^{i})$ 的上同调都集中在区间 $[z',z'+n']$ 中，故对任一 $i<0$ 与 $j\notin[-z'-n',-z'+\dim M]$ 都有 $\Hom^{j}(G(\cL^{i}),\cF)=0$。（这里用到：范畴 $\coh(M)$ 的同调维数等于 $\dim M$。）因此，对这些同样的 $j$，对 $M$ 上任一层 $\cF$ 都有 $H^{j}(F(\cF))=0$。于是函子 $F$ 有界。
\end{proof}

\begin{remark}{3.2.4}
必要时在导出范畴中平移 $F$ 之后，我们自现在起并贯穿本章假设：对 $M$ 上的任一层 $\cF$，上同调 $H^{i}(F(\cF))$ 仅在 $i\in[-a,0]$ 时非零，其中 $a$ 是取定的正整数。
\end{remark}

% [Tutor Note: 本块止于注 3.2.4；“3.3. Construction of the object representing a fully faithful functor”标题及其后内容归 ch03b，此处不译。]
